 
% new file names from date 20171119: 
% block0_171119.tex 
\renewcommand{\thisfile}{block0\_190328}
 
\cleardoublepage
%\bigskip
% \setcounter{chapter}{-1}
\setcounter{section}{-1}
\mysection{Introduction}
\AYcopyright


\mysubsection{On the Subject} \label{subsec:on-the-sub}
{\em Derived categories} were introduced by A. Grothendieck and J.-L.Verdier 
around 1960, and were first published in the book \cite{RD} by R. 
Hartshorne. The basic idea was as follows. They had realized that the derived 
functors of classical homological algebra, namely the functors 
$\mrm{R}^q F, \, \mrm{L}_q F : \cat{M} \to \cat{N}$ 
derived from an additive functor 
$F : \cat{M} \to \cat{N}$
between abelian categories, are too limited to allow several rather natural 
manipulations. Perhaps the most important operation that was lacking was the 
composition of derived functors; the best approximation of it was a spectral 
sequence. 

The solution to the problem was to invent a new category, starting from a given 
abelian category $\cat{M}$. The objects of this new category are the complexes 
of objects of $\cat{M}$. These are the same complexes that play an auxiliary 
role in classical homological algebra, as resolutions of objects of $\cat{M}$. 
The complexes form a category $\dcat{C}(\cat{M})$, but this category is not 
sufficiently intricate to carry in it the information of derived 
functors. So it must be modified. 

A morphism $\phi : M \to N$ in $\dcat{C}(\cat{M})$
is called a {\em quasi-isomorphism}
if in each degree $q$ the cohomology morphism 
$\opn{H}^q(\phi) : \opn{H}^q(M) \to \opn{H}^q(N)$
in $\cat{M}$ is an isomorphism. The modification that is needed is to make the 
quasi-isomorphisms invertible. This is done by a formal 
localization procedure, and the resulting category (with the same objects as 
$\dcat{C}(\cat{M})$) is the derived category $\dcat{D}(\cat{M})$.
There is a functor 
$\opn{Q} : \dcat{C}(\cat{M}) \to \dcat{D}(\cat{M})$,
which is the identity on objects, and it has a universal property (it is 
initial among the functors that send the quasi-isomorphisms to isomorphisms).  
A theorem (analogous to Ore localization in noncommutative ring theory) says 
that every morphism $\th$ in $\dcat{D}(\cat{M})$ can be written as a simple 
left 
or right fraction:
\begin{equation} \label{eqn:4994}
\th = \opn{Q}(\psi_0)^{-1} \circ \opn{Q}(\phi_0) = 
\opn{Q}(\phi_1) \circ \opn{Q}(\psi_1)^{-1} ,
\end{equation}
where $\phi_i$ and $\psi_i$ are morphisms in $\dcat{C}(\cat{M})$, and 
$\psi_i$ are quasi-isomorphisms. 

The cohomology functors 
$\opn{H}^q : \dcat{D}(\cat{M}) \to \cat{M}$,
for all $q \in \Z$, are still defined. 
It turns out that the functor $\cat{M} \to \dcat{D}(\cat{M})$,
that sends an object $M$ to the complex $M$ concentrated in degree $0$, is 
fully faithful. 

The next step is to say what is a left or a right derived functor of an 
additive functor $F : \cat{M} \to \cat{N}$. The functor $F$ can be extended 
in an obvious manner to a functor on complexes
$F : \dcat{C}(\cat{M}) \to \dcat{C}(\cat{N})$.
A {\em right derived functor} of $F$ is a functor 
\begin{equation} \label{eqn:4990}
\mrm{R} F : \dcat{D}(\cat{M}) \to \dcat{D}(\cat{N}) ,
\end{equation}
together with a morphism of functors 
$\eta^{\mrm{R}} : \opn{Q}_{\cat{N}} \circ \, F \to \mrm{R} F \circ 
\opn{Q}_{\cat{M}}$.
The pair $(\mrm{R} F, \eta^{\mrm{R}})$ has to be {\em initial among all such 
pairs}. The uniqueness of such a functor $\mrm{R} F$, up to a unique 
isomorphism, is relatively easy to prove (using the language of 
$2$-categories). As for existence of $\mrm{R} F$, it relies on the 
existence of suitable resolutions (like the injective resolutions in the 
classical situation). If these resolutions exist, 
and if the original functor $F$ is left exact,
then there is a canonical isomorphism of functors 
\begin{equation} \label{eqn:4991}
\mrm{R}^q F \cong \opn{H}^q \circ \, \mrm{R} F : \cat{M} \to \cat{N} 
\end{equation}
for every $q \geq 0$. 

The left derived functor 
\begin{equation} \label{eqn:4992}
\mrm{L} F : \dcat{D}(\cat{M}) \to \dcat{D}(\cat{N}) 
\end{equation}
is defined similarly. When suitable resolutions exist, and when 
$F$ is right exact, there is a canonical isomorphism of functors
\begin{equation} \label{eqn:4993}
\mrm{L}_q F \cong \opn{H}^{-q} \circ \, \mrm{L} F : \cat{M} \to \cat{N} 
\end{equation}
for every $q \geq 0$. 

There are several variations: $F$ could be a contravariant additive 
functor; or it could be an additive bifunctor, contravariant in one or two of 
its arguments. In all these situations the derived (bi)functors 
$\mrm{R} F$ and $\mrm{L} F$ can be defined. 

The derived category $\dcat{D}(\cat{M})$ is additive, but it is not abelian. 
The notion of short exact sequence (in $\cat{M}$ and in $\dcat{C}(\cat{M})$) is 
replaced by that of {\em distinguished 
triangle}, and thus $\dcat{D}(\cat{M})$ is a {\em triangulated category}. 
The derived functors $\mrm{R} F$ and $\mrm{L} F$ are {\em triangulated 
functors}, which means that they send distinguished triangles in 
$\dcat{D}(\cat{M})$ to distinguished triangles in $\dcat{D}(\cat{N})$. 

Already in classical homological algebra we are interested in the {\em 
bifunctors} $\opn{Hom}(-, -)$ and $(- \ot -)$. These bifunctors can also be 
derived. To simplify matters, let's assume that $A$ is a commutative ring, and 
$\cat{M} = \cat{N} = \cat{Mod} A$, the category of $A$-modules. 
We then have bifunctors 
\[ \opn{Hom}_A(-, -) \, : \, (\cat{Mod} A)^{\mrm{op}} \, \times \, \cat{Mod} A
\, \to \, \cat{Mod} A \]
and 
\[ (- \ot_A -) \, : \, \cat{Mod} A \, \times \, \cat{Mod} A \, \to \, 
\cat{Mod} A \, , \]
where the superscript ``op'' denotes the opposite category, that encodes the
contravariance in the first argument of Hom. 
In this situation all resolutions exist, and we have the right derived 
bifunctor 
\begin{equation} \label{eqn:4995}
\opn{RHom}_A(-, -) \, : \, \dcat{D}(\cat{Mod} A)^{\mrm{op}} \, \times \, 
\dcat{D}(\cat{Mod} A) \, \to \, \dcat{D}(\cat{Mod} A) 
\end{equation}
and the left derived bifunctor 
\begin{equation} \label{eqn:4996}
(- \ot^{\mrm{L}}_{A} -) \, : \, \dcat{D}(\cat{Mod} A) \, \times \, 
\dcat{D}(\cat{Mod} A) \, \to \, \dcat{D}(\cat{Mod} A) \, .
\end{equation}
The compatibility with the classical derived bifunctors is this: there are 
canonical isomorphisms 
\begin{equation} \label{eqn:4997}
\opn{Ext}^q_A(M, N) \, \cong \,
\opn{H}^q \bigl( \opn{RHom}_A(M, N) \bigr) 
\end{equation}
and
\begin{equation} \label{eqn:4998}
\opn{Tor}_q^A(M, N) \, \cong  \, \opn{H}^{-q} (M \ot^{\mrm{L}}_{A} N) 
\end{equation}
for all $M, N \in \cat{Mod} A$ and $q \geq 0$.  

This is what derived categories and derived functors are. As to what can be 
done with them, here are some of the things we will explore in our book:
\begin{itemize}
\item {\em Dualizing complexes} and {\em residue complexes} over noetherian 
commutative rings. Besides the original treatment from \cite{RD}, that we 
present in detail here, we also include {\em Van den Bergh rigidity} in the 
commutative setting, that gives rise to {\em rigid residue complexes}. 

\item {\em Perfect DG modules} and {\em tilting DG bimodules}  over 
noncommutative DG rings, and a few variants of {\em derived 
Morita Theory}, including the {\em Rickard-Keller Theorem}. 

\item {\em Derived torsion} and {\em balanced dualizing complexes}
over connected graded NC rings, and {\em rigid 
dualizing complexes} over NC rings, including a full proof of the {\em Van den 
Bergh Existence Theorem} for NC rigid dualizing complexes. 
\end{itemize}

A topic that is beyond the scope of this book, but of which we 
provide an outline here, is:
\begin{itemize}
\item The {\em rigid approach to Grothendieck Duality} on 
noetherian schemes and Deligne-Mumford stacks. 
\end{itemize}

Derived categories have important roles in several areas of mathematics; below 
is a partial list. We will not be able to talk about any of these 
topics in our book, so instead we give some references alongside each 
topic. 
\begin{itemize}
\item[$\triangleright$] $\mcal{D}$-modules, perverse sheaves, and 
representations of algebraic groups and Lie algebras.  
See \cite{BBD} and \cite{Bor}. 
More recently the focus in this area is on the {\em Geometric Langlands 
Correspondence}, that can only be stated in terms of derived categories (see 
the survey \cite{Gai}).

\item[$\triangleright$] 
Algebraic analysis, including differential, 
microdifferential and DQ modules (see \cite{Ka}, \cite{SKK}, \cite{KaSc3}) 
and microlocal sheaf theory (see \cite{KaSc1}), with its application to 
symplectic topology (see \cite{Ta}, \cite{NaZa}). 

\item[$\triangleright$] Representations of finite groups and quivers, 
including {\em cluster algebras} and the {\em Brou\'e Conjecture}. 
See \cite{Hap}, \cite{Kel2}, \cite{CrRo}. 

\item[$\triangleright$] Birational algebraic geometry. This includes {\em 
Fourier-Mukai transforms} and {\em semi-orthogonal decompositions}. 
See the surveys \cite{HiVdB} and \cite{MacSt}, and the book \cite{Huy}. 

\item[$\triangleright$] Homological mirror symmetry. It relates the derived 
category of coherent sheaves on a complex algebraic variety $X$ to the 
{\em derived Fukaya category} of the mirror partner $Y$, which is a symplectic 
manifold. See Remark \ref{rem:4960} and the online reference \cite{HMS}.  

\item[$\triangleright$] Derived algebraic geometry. Here not only the category 
of sheaves is derived, but also the underlying geometric objects (schemes or 
stacks). See Example \ref{exa:4205}, Remark \ref{rem:4970}, and the references 
\cite{Lur} and \cite{To2}.  
\end{itemize}


\mysubsection{A Motivating Discussion: Duality} \label{subsec:2160}
Let us now approach derived categories from another perspective,
very different from the one taken in the previous subsection, by 
considering the idea of {\em duality in algebra}. 

We begin with something elementary: linear algebra. 
Take a field $\K$. Given a $\K$-module $M$ (i.e.\  a vector space), let
$D(M) := \opn{Hom}_{\K}(M, \K)$,
the dual module. There is a canonical homomorphism 
\begin{equation} \label{eqn:3250}
\opn{ev}_M : M \to D(D(M)) ,
\end{equation}
called {\em Hom-evaluation}, whose formula is 
$\opn{ev}_M(m)(\phi) := \phi(m)$ for 
$m \in M$ and $\phi \in D(M)$. 
If $M$ is finitely generated then $\opn{ev}_M$ is an isomorphism (actually this 
is ``if and only if''). 

To formalize this situation, let $\cat{Mod} \K$ denote the category of
$\K$-modules. Then 
$D : \cat{Mod} \K \to \cat{Mod} \K$
is a contravariant functor, and 
$\opn{ev} : \opn{Id} \to D \circ D$
is a morphism of functors (i.e.\ a natural transformation). Here $\opn{Id}$ is 
the identity functor of $\cat{Mod} \K$. 

Now let us replace $\K$ by some nonzero commutative ring $A$. Again we can
define a contravariant functor 
\begin{equation} \label{eqn:3241}
D : \cat{Mod} A \to \cat{Mod} A , \quad
D(M) := \opn{Hom}_A(M, A) ,
\end{equation}
and a morphism of functors $\opn{ev} : \opn{Id} \to D \circ D$.
It is easy to see that $\opn{ev}_M : M \to D(D(M))$ is an isomorphism if $M$
is a finitely generated free $A$-module. 
Of course we can't expect reflexivity (i.e.\ $\opn{ev}_M$ being an isomorphism) 
if $M$ is not finitely generated; but what about a finitely generated module 
that is not free?

In order to understand this better, let us concentrate on the ring $A = \Z$.
Since $\Z$-modules are just abelian groups, the category $\cat{Mod} \Z$ is 
often denoted by $\cat{Ab}$. 
Let $\cat{Ab}_{\mrm{f}}$ be the full subcategory of finitely generated abelian 
groups. Every finitely generated abelian group is of the form 
$M \cong T \oplus F$, with $T$ finite and $F$ free. (The letters  ``T'' and 
``F''  stand for ``torsion'' and ``free'' respectively.) It is important to 
note 
that this is {\em not a canonical 
isomorphism}. There is a canonical short exact sequence
\begin{equation} \label{eqn:1000}
0 \to T \xar{\phi} M \xar{\psi} F \to 0 
\end{equation}
in $\cat{Ab}_{\mrm{f}}$, but the decomposition $M \cong T \oplus F$ comes from 
{\em choosing a splitting} $\sigma : F \to M$ of this sequence. 

\begin{exer}
Prove that the exact sequence (\ref{eqn:1000}) is functorial; 
namely there are functors 
$T, F : \cat{Ab}_{\mrm{f}} \to \cat{Ab}_{\mrm{f}}$,
and natural transformations
$\phi : T \to \opn{Id}$ and  
$\psi :  \opn{Id} \to F$, such that for each $M \in \cat{Ab}_{\mrm{f}}$ 
the group $T(M)$ is finite, the group $F(M)$ is free, and 
the sequence of homomorphisms
\begin{equation} \label{eqn:1002}
0 \to T(M) \xar{\phi_M} M  \xar{\psi_M} F(M) \to  0
\end{equation}
is exact. 

Next, prove that there does not exist a {\em functorial decomposition} of a
finitely generated abelian group into a free part and a finite part. 
Namely, there is no natural transformation
$\sigma : F \to \opn{Id}$, such that for every $M$ the homomorphism 
$\sigma_M : F(M) \to M$ 
splits the sequence (\ref{eqn:1002}). (Hint: find a counterexample.)
\end{exer}

We know that for a free finitely generated abelian group $F$ there is 
reflexivity, i.e.\ $\opn{ev}_F : F \to D(D(F))$ is an isomorphism. But for a
finite abelian group $T$ we have
$D(T) = \opn{Hom}_{\Z}(T, \Z) = 0$.
Thus, for a group $M \in \cat{Ab}_{\mrm{f}}$  with a nonzero torsion subgroup
$T$, reflexivity fails: $\opn{ev}_M : M \to D(D(M))$ is not an isomorphism.

On the other hand, for an abelian group $M$ we can define another sort of dual:
$D'(M) := \opn{Hom}_{\Z}(M, \mbb{Q} / \Z)$.
There is a morphism of functors 
$\opn{ev}' : \opn{Id} \to D' \circ D'$. 
For a finite abelian group $T$ the homomorphism 
$\opn{ev}'_T : T \to D'(D'(T))$ is an isomorphism; this can be seen by 
decomposing $T$ into cyclic groups, and for a finite cyclic group it is clear. 
So $D'$ is a duality for finite abelian groups.
(We may view the abelian group $\mbb{Q} / \Z$ as the group of roots of $1$ in
$\mbb{C}$, via the exponential function; and then $D'$ becomes {\em 
Pontryagin Duality}.)

But for a finitely generated free abelian group $F$ we get
$D'(D'(F)) = \what{F}$, the profinite completion of $F$. So once more this is
not a good duality for all finitely generated abelian groups.

This is where the {\em derived category} enters. 
For every commutative ring $A$ there is the derived category
$\dcat{D}(\cat{Mod} A)$. Here is a very quick explanation of it,
in concrete terms -- as opposed to the abstract point of view taken in the 
previous subsection. 

Recall that a {\em complex} of $A$-modules is a diagram
\begin{equation} \label{eqn:1010}
M = \bigl(  \cdots \to M^{-1} \xar{\d_M^{-1}} M^{0} 
\xar{\d_M^0} M^{1} \to \cdots \bigr) 
\end{equation}
in the category $\cat{Mod} A$.
Namely the $M^i$ are $A$-modules, and the $\d_M^{i}$ are homomorphisms.
The condition is that $\d_M^{i + 1} \circ \d_M^{i} = 0$.
We sometimes write 
$M = \{ M^i \}_{i \in \Z}$. 
The collection $\d_M = \{ \d_M^{i} \}_{i \in \Z}$ is called the 
{\em differential} of $M$.

Given a second complex 
\[ N = \bigl(  \cdots \to N^{-1} \xar{\d_N^{-1}} N^{0} 
\xar{\d_N^0} N^{1} \to \cdots \bigr) , \]
a {\em homomorphism of complexes} 
$\phi^{} : M^{} \to N^{}$
is a collection 
$\phi = \{ \phi^{i} \}_{i \in \Z}$ 
of homomorphisms 
$\phi^i : M^i \to N^i$ 
in $\cat{Mod} A$ satisfying 
$\phi^{i + 1} \circ \d_M^{i} =  \d_N^{i} \circ \phi^i$.
The resulting category is denoted by $\dcat{C}(\cat{Mod} A)$. 

The $i$-th {\em cohomology} of the complex $M$ is 
\begin{equation} \label{eqn:5000}
\opn{H}^i(M^{}) := 
\opn{Ker}(\d_M^{i}) \, / \, \opn{Im}(\d_M^{i - 1}) \in \cat{Mod} A . 
\end{equation}
A homomorphism  $\phi^{} : M^{} \to N^{}$
in $\dcat{C}(\cat{Mod} A)$ induces homomorphisms 
$\opn{H}^i(\phi^{}) : \opn{H}^i(M^{}) \to \opn{H}^i(N^{})$ 
in $\cat{Mod} A$. We call $\phi^{}$ a {\em quasi-isomorphism} if all the 
homomorphisms $\opn{H}^i(\phi^{})$ are isomorphisms. 

The derived category $\dcat{D}(\cat{Mod} A)$ is the localization of 
$\dcat{C}(\cat{Mod} A)$ with respect to the quasi-isomorphisms. 
This means that $\dcat{D}(\cat{Mod} A)$ has the same objects as 
$\dcat{C}(\cat{Mod} A)$. There is a functor 
\begin{equation} \label{eqn:5001}
\opn{Q} : \dcat{C}(\cat{Mod} A) \to \dcat{D}(\cat{Mod} A)
\end{equation}
which is the identity on objects, it sends quasi-isomorphisms to 
isomorphisms, and it is universal for this property. 

A single $A$-module $M^0$ can be viewed as a complex $M$ concentrated in degree 
$0$:
\[ M^{} = \bigl(  \cdots \to 0 \xar{0} M^0 \xar{0} 0 \to \cdots \bigr) . \]
This turns out to be a fully faithful embedding 
\begin{equation} \label{eqn:1004}
\cat{Mod} A \to \dcat{D}(\cat{Mod} A) .
\end{equation}
The essential image of this embedding is the full subcategory 
of $\dcat{D}(\cat{Mod} A)$ on the complexes $M^{}$ whose cohomology is 
concentrated in degree $0$. In this way we have {\em enlarged} the category of 
$A$-modules. All this is explained in Sections 
\ref{sec:loc-cats}-\ref{sec:der-cat} of the book. 

Here is a very important kind of quasi-isomorphism. Suppose $M$ is an 
$A$-module and 
\begin{equation} \label{eqn:1020}
\cdots \to P^{-2} \xar{\d_P^{-2}}  P^{-1} \xar{\d_P^{-1}}  P^{0} 
\xar{\rho} M \to 0 
\end{equation}
is a projective resolution of it.
We can view $M$ as a complex concentrated in degree $0$, by the embedding 
(\ref{eqn:1004}). Define the complex 
\begin{equation} \label{eqn:3246}
P := \bigl(  \cdots \to P^{-2} \xar{\d_P^{-2}} 
P^{-1} \xar{\d_P^{-1}} P^0 \to 0  \to \cdots \bigr) ,
\end{equation}
concentrated in nonpositive degrees. 
Then $\rho$ becomes a morphism of complexes 
$\rho : P \to  M$.
The exactness of the sequence (\ref{eqn:1020}) says that $\rho$ is actually a 
quasi-iso\-morphism. 
Thus $\opn{Q}(\rho) : P^{} \to M$
is an isomorphism in $\dcat{D}(\cat{Mod} A)$.

Let us fix a complex $R \in \dcat{C}(\cat{Mod} A)$.
For every complex $M \in \dcat{C}(\cat{Mod} A)$ we can form the complex
\[ D(M) := \opn{Hom}_{A}(M, R) \in \dcat{C}(\cat{Mod} A) . \]
This is the usual Hom complex (we recall it in 
Subsection \ref{subsec:complexes}). 
As $M$ changes we get a contravariant functor 
\[ D : \dcat{C}(\cat{Mod} A) \to \dcat{C}(\cat{Mod} A) . \]
The functor $D$ has a {\em contravariant right derived functor} 
\begin{equation} \label{eqn:5002}
\mrm{R} D : \dcat{D}(\cat{Mod} A) \to \dcat{D}(\cat{Mod} A) .
\end{equation}
If $P$ is a bounded above complex of projective modules (like in formula 
(\ref{eqn:3246})), or more generally a {\em K-projective complex}
(see Subsection \ref{subsec:K-proj}), then there is a canonical isomorphism
\begin{equation} \label{eqn:3242}
\mrm{R} D(P) \cong D(P) = \opn{Hom}_A(P, R) . 
\end{equation}
Every complex $M$ admits a K-projective resolution
$\rho : P \to M$, and this allows us to calculate $\mrm{R} D(M)$.
Indeed, because the morphism 
$\opn{Q}(\rho) : P \to M$ 
is an isomorphism in 
$\dcat{D}(\cat{Mod} A)$, it follows that 
$\mrm{R} D(\opn{Q}(\rho)) : \mrm{R} D(M) \to \mrm{R} D(P)$
is an isomorphism in $\dcat{D}(\cat{Mod} A)$. 
And the complex $\mrm{R} D(P)$ is known by the canonical isomorphism 
(\ref{eqn:3242}). All this is explained in Sections \ref{sec:der-funcs}, 
\ref{sec:resol} and \ref{sec:exist-resol} of the book. 

It turns out that there is a canonical morphism 
\begin{equation} \label{eqn:3247}
\opn{ev}^{\mrm{R}} : \opn{Id} \to \mrm{R} D \circ \mrm{R} D
\end{equation}
of functors from $\dcat{D}(\cat{Mod} A)$ to itself, called {\em derived 
Hom-evaluation}. See Subsection 
\ref{subsec:du-cplxs}. 

Let us now return to the ring $A = \Z$ and the complex $R = \Z$. 
So the functor $D$ is the same one we had in (\ref{eqn:3241}).
Given a finitely generated abelian group $M$, we want to calculate the 
complexes $\mrm{R} D(M)$ and $\mrm{R} D(\mrm{R} D(M))$, and the morphism 
\begin{equation} \label{eqn:3248}
\opn{ev}^{\mrm{R}}_M : M \to \mrm{R} D(\mrm{R} D(M)) 
\end{equation}
in $\dcat{D}(\cat{Mod} A)$. As explained above, for this we choose a 
projective resolution $\rho : P \to M$,
and then we calculate  the complexes $\mrm{R} D(P)$ and 
$\mrm{R} D(\mrm{R} D(P))$, and the morphism $\opn{ev}^{\mrm{R}}_P$. 
For convenience we choose a projective resolution $P$ of this shape:
\[ \begin{aligned}
& P^{} = 
\bigl(  \cdots \to 0 \to P^{-1} \xar{\d_P^{-1}} P^0 \to 0 \to \cdots \bigr)
\\ & \quad \ 
= \bigl( \cdots \to 0 \xar{}  \Z^{r_1} \xar{\bsym{a} \cd (-)} \Z^{r_0} \xar{} 0 
\cdots 
\bigr) \, , 
\end{aligned} \]
where $r_0, r_1 \in \N$ and $\bsym{a}$ is a matrix of integers. 
The complex $\mrm{R} D(P)$ is this: 
\[ \mrm{R} D(P) \cong D(P) = \opn{Hom}_{\Z}(P^{}, \Z) =
\bigl( \cdots \to 0 \xar{}  \Z^{r_0} \xar{\bsym{a}^{\mrm{t}} \cd (-)}  
\Z^{r_1} \xar{} 0 \cdots \bigr) , \]
a complex of free modules concentrated in degrees $0$ and $1$, with the 
transpose matrix $\bsym{a}^{\mrm{t}}$ as its differential. 

Because $\mrm{R} D(P^{}) \cong D(P^{})$ is itself a bounded complex of 
free modules, its derived dual is 
\begin{equation} \label{eqn:5003}
\mrm{R} D(\mrm{R} D(P^{})) \cong D(D(P^{})) = 
\opn{Hom}_{\Z} \bigl( \opn{Hom}_{\Z}(P^{}, \Z), \Z \bigr) .
\end{equation}
Under the isomorphism (\ref{eqn:5003}), the derived Hom-evaluation morphism 
$\opn{ev}^{\mrm{R}}_P$ in this case is 
just the naive Hom-evaluation homomorphism 
$\opn{ev}_P : P^{} \to D(D(P^{}))$ 
in $\dcat{C}(\cat{Mod} \Z)$ from (\ref{eqn:3250}); see Exercise 
\ref{exer:3245}. Because $P^0$ and $P^{-1}$ are finite rank free modules, it 
follows that $\opn{ev}_P$ is an isomorphism in $\dcat{C}(\cat{Mod} \Z)$. 
Therefore the morphism $\opn{ev}^{\mrm{R}}_M$ in $\dcat{D}(\cat{Mod} \Z)$
is an isomorphism. We see that  {\em $\mrm{R} D$ is a duality that holds for 
all finitely generated $\Z$-modules} $M$~!

Actually, much more is true. Let us denote by 
$\dcat{D}_{\mrm{f}}(\cat{Mod} \Z)$
the full subcategory of $\dcat{D}(\cat{Mod} \Z)$ on the complexes $M$
such that $\opn{H}^i(M)$ is finitely generated for all $i$. Then, according to 
Theorem \ref{thm:2155}, $\opn{ev}^{\mrm{R}}_M$ is an isomorphism for every 
$M \in \dcat{D}_{\mrm{f}}(\cat{Mod} \Z)$. 
It follows that 
\begin{equation} \label{eqn:5004}
\mrm{R} D : \dcat{D}_{\mrm{f}}(\cat{Mod} \Z) \to 
\dcat{D}_{\mrm{f}}(\cat{Mod} \Z)
\end{equation}
is a duality (a contravariant equivalence). This is the celebrated {\em 
Grothendieck Duality}. 

Here is the connection between the derived duality $\mrm{R} D$ and the 
classical dualities $D$ and $D'$. Take a finitely generated abelian group 
$M$, with short exact sequence (\ref{eqn:1000}). There are
canonical isomorphisms
\[ \mrm{H}^0 (\mrm{R} D(M)) \cong \opn{Ext}^0_{\Z}(M, \Z) \cong 
\opn{Hom}_{\Z}(M, \Z) \cong \opn{Hom}_{\Z}(F, \Z) = D(F)  \]
and 
\[ \mrm{H}^1 (\mrm{R} D(M)) \cong \opn{Ext}^1_{\Z}(M, \Z) 
\cong \opn{Ext}^1_{\Z}(T, \Z) \cong D'(T)  .  \]
The cohomologies $\mrm{H}^i (\mrm{R} D(M))$ vanish for $i \neq 0, 1$. 
We see that if $M$ is neither free nor finite, then  
$\mrm{H}^0 (\mrm{R} D(M))$ and $\mrm{H}^1 (\mrm{R} D(M))$ are both nonzero; 
so that the complex $\mrm{R} D(M)$ is not isomorphic
in $\dcat{D}(\cat{Mod} \Z)$ to an object of 
$\cat{Mod} \Z$, under the embedding (\ref{eqn:1004}). 

Grothendieck Duality holds for many noetherian commutative rings $A$. 
A sufficient condition is that $A$ is a finitely generated ring over a 
regular noetherian ring $\K$ (e.g.\ $\K = \Z$ or a field). 
A complex $R \in \dcat{D}(\cat{Mod} A)$ for which the contravariant functor 
\begin{equation} \label{eqn:5010}
\mrm{R} D = \opn{RHom}_{A}(-, R) : \dcat{D}_{\mrm{f}}(\cat{Mod} A) \to 
\dcat{D}_{\mrm{f}}(\cat{Mod} A)
\end{equation}
is an equivalence is called a  {\em dualizing complex}. 
(This is not quite accurate -- see Definition \ref{dfn:2155} for the precise 
technical conditions on $R$.) 
A dualizing complex $R$ over $A$ is unique (up to a degree translation and 
tensoring with an invertible module). 
See Theorems \ref{thm:2155}, \ref{thm:2170} and  \ref{thm:2175}. 

Interestingly, the structure of the dualizing complex $R$ depends on the
geometry of the ring $A$ (i.e.\ of the affine scheme $\opn{Spec}(A)$). 
If $A$ is a regular ring (like $\Z$) then $R = A$ is dualizing. 
If $A$ is a Cohen-Macaulay ring (and $\opn{Spec}(A)$ is connected) then $R$ is 
a single $A$-module (up to a shift in degrees). But if $A$ is a more 
singular ring, then $R$ must live in several degrees, as the next example 
shows. 

\begin{exa} \label{exa:2252}
Consider the affine algebraic variety 
$X \subseteq \mbf{A}^3_{\mbb{R}}$
which is the union of a plane and a line that meet at a point, with 
coordinate ring
\[ A = \mbb{R}[t_1, t_2, t_3] / (t_3 \cd t_1, \, t_3 \cd t_2) . \]
See figure \ref{fig:1}. 
A dualizing complex $R$ over $A$ must live in two adjacent degrees;
namely there is some $i$ such that both 
$\mrm{H}^i(R)$ and $\mrm{H}^{i+1}(R)$ are nonzero. 
This calculation is worked out in full in Example \ref{exa:2245}.
\end{exa}

\begin{figure}[!htb]
\centering
\includegraphics[scale=0.46]{fig1.jpg}
\caption{An algebraic variety $X$ that is connected but not equidimensional, 
and hence not Cohen-Macaulay.} 
\label{fig:1}
\end{figure}

One can also talk about dualizing complexes over {\em noncommutative rings}. 
We will do this in Sections \ref{sec:BDC} and \ref{sec:rigid-DC-NC}.





\mysubsection{On the Book}

This book develops the theory of derived categories, starting from the 
foundations, and going all the way to applications in commutative and 
noncommutative algebra. 
The emphasis is on explicit constructions (with examples), as opposed to 
axiomatics. The most abstract concept we use is probably that of an abelian 
category (which seems indispensable). 

A special feature of this book is that most of the theory deals with 
the category $\dcat{C}(A, \cat{M})$ of {\em DG $A$-modules in $\cat{M}$}, 
where $A$ is a DG  ring and $\cat{M}$ is an 
abelian category. Here ``DG'' is short for ``differential graded'', and our 
DG rings are more commonly known as unital associative DG algebras.
See Subsections \ref{subsec:DG-rng-mod} and 
\ref{subsec:DGModinM} for the definitions. 
The notion $\dcat{C}(A, \cat{M})$ covers most important examples that 
arise in algebra and geometry: 
\begin{itemize}
\item The category $\dcat{C}(A)$ of DG $A$-modules, for any DG ring 
$A$. This includes unbounded complexes of modules over an ordinary ring $A$. 

\item The category $\dcat{C}(\cat{M})$ of unbounded complexes in any 
abelian category $\cat{M}$. This includes 
$\cat{M} = \cat{Mod} \AA$, the category of sheaves of $\AA$-modules on a ringed 
space $(X, \AA)$. 
\end{itemize}

The category $\dcat{C}(A, \cat{M})$ is a {\em DG category}, and 
its DG structure determines the {\em homotopy category} 
$\dcat{K}(A, \cat{M})$ with its {\em triangulated structure}. 
We prove that every {\em DG functor} 
$F : \dcat{C}(A, \cat{M}) \to \dcat{C}(B, \cat{N})$
induces a {\em triangulated functor} 
\begin{equation} \label{eqn:3205}
F : \dcat{K}(A, \cat{M}) \to \dcat{K}(B, \cat{N})
\end{equation}
between the homotopy categories. 

We can now reveal that in the previous subsections we were a bit imprecise (for 
the sake of simplifying the exposition): what we referred to there as  
$\dcat{C}(\cat{M})$ was actually the {\em strict subcategory}
$\dcat{C}_{\mrm{str}}(\cat{M})$, whose morphisms are the degree $0$ cocycles in 
the DG category $\dcat{C}(\cat{M})$. For the same reason the homotopy 
category $\dcat{K}(\cat{M})$ was suppressed there. 

The {\em derived category} $\dcat{D}(A, \cat{M})$ is obtained 
from $\dcat{K}(A, \cat{M})$ by inverting the quasi-isomorphisms.
If suitable resolutions exist, the triangulated functor (\ref{eqn:3205}) 
can be derived, on the left or on the right, to yield triangulated functors
\begin{equation} \label{eqn:5011}
\mrm{L} F, \, \mrm{R} F : 
\dcat{D}(A, \cat{M}) \to \dcat{D}(B, \cat{N}) .
\end{equation}
We prove existence of {\em K-injective}, {\em K-projective} and {\em K-flat} 
resolutions in $\dcat{K}(A, \cat{M})$ in several important contexts,
and explain their roles in deriving functors and in presenting morphisms
in $\dcat{D}(A, \cat{M})$.

In the last six sections of the book we discuss a few key applications of 
derived categories to commutative and noncommutative algebra. 

There is a section-by-section synopsis of the book in subsection 
\ref{subsec:synopsis} of the Introduction. 

The book is based on notes for advanced courses given at Ben Gurion 
University, in the academic years 2011-12, 2015-16 and 2016-17. 
The main sources for Sections 1-12 of the book are \cite{RD} and \cite{KaSc1}; 
but the DG theory component is absent from those earlier texts, and is pretty 
much our own interpretation of folklore results. The material covered in 
Sections 13-18 is adapted from research papers. 


\mysubsection{Synopsis of the Book} \label{subsec:synopsis}
Here is a section-by-section description of the material in the book. 

\medskip \noindent 
{\bf Sections \ref{sec:basic}-\ref{sec:abelian}.} These sections are pretty 
much a review of the standard material on categories and functors, especially 
abelian categories and additive functors, that is needed for the book. 

There is one new topic here -- it is our {\em sheaf tricks}, in Subsection 
\ref{subsec:sheaf-tricks}, that are a ``cheap substitute'' for the 
Freyd-Mitchell Theorem, in the sense that they facilitate proofs of various 
results in abstract abelian categories, yet they are themselves easy to prove 
(and we provide these proofs). 

A reader who is familiar with this material can skip these sections; yet we 
do recommend looking at our notational conventions, that are spelled out in 
Conventions \ref{conv:2490} and \ref{conv:3170}. 

\medskip \noindent 
{\bf Section \ref{sec:DG-algebra}.} A good understanding of {\em DG algebra} 
(``DG'' is short for ``differential graded'') is essential in our approach to 
derived categories. By DG algebra we mean DG rings, DG modules, DG categories 
and DG functors. There do not exist (to our knowledge) detailed textbook 
references for DG algebra. Therefore we have included a lot of basic material 
in this section.

We work over a fixed nonzero commutative base ring $\K$ (e.g.\ a field or the 
ring of integers $\Z$). All linear operations (rings, categories, functors, 
etc.) are assumed to be $\K$-linear. A {\em DG module} 
is a module $M$ with a direct sum decomposition 
$M = \boplus_{i \in \Z} M^i$ into submodules, with a differential $\d_M$ of 
degree $1$ satisfying $\d_M \circ \d_M = 0$. The grading on $M$ is called  
{\em cohomological grading}. 
Tensor products of DG modules have {\em signed braiding isomorphisms}: for 
DG modules $M$ and $N$, and for homogeneous elements 
$m \in M^i$ and $n \in N^j$, we define 
\[ \opn{br}_{M, N}(m \ot n) := (-1)^{i \cd j} \cd n \ot m \in N \ot_{\K} M . \]
This braiding is usually referred to as the {\em Koszul sign rule}. 

A {\em DG category} is a $\K$-linear
category $\cat{C}$ in which the morphism sets 
$\opn{Hom}_{\cat{C}}(M, N)$ are endowed with DG module structures, 
and the composition functions are DG bilinear.  
There are two important sources of DG categories: the category 
$\dcat{C}(\cat{M})$ of complexes in a $\K$-linear abelian category $\cat{M}$,
and the category of DG modules over a central DG $\K$-ring $A$
(traditionally called a unital associative DG $\K$-algebra).
Since we want to consider both setups, but we wish to avoid repetition, we 
have devised a new concept that combines both: the DG category 
$\dcat{C}(A, \cat{M})$ of {\em DG $A$-modules in $\cat{M}$}. 
By definition, a DG $A$-module in $\cat{M}$ is a pair $(M, f)$, consisting of a 
complex $M \in \dcat{C}(\cat{M})$, together with a DG $\K$-ring homomorphism 
$f : A \to \opn{End}_{\cat{M}}(M)$.
The morphisms in $\dcat{C}(A, \cat{M})$ are the morphism in 
$\dcat{C}(\cat{M})$ that respect the action of $A$. 
When $A = \K$ we are in the case 
$\dcat{C}(A, \cat{M}) = \dcat{C}(\cat{M})$, and when 
$\cat{M} = \cat{Mod} \K$ we are in the case 
$\dcat{C}(A, \cat{M}) = \dcat{C}(A)$. 

A morphism $\phi : M \to N$ in $\dcat{C}(A, \cat{M})$
is called a {\em strict morphism} if it is a degree $0$ cocycle. 
The {\em strict subcategory} $\dcat{C}_{\mrm{str}}(A, \cat{M})$
is the subcategory of $\dcat{C}(A, \cat{M})$ on all the objects, 
and its morphisms are the strict morphisms. The strict category is abelian. 

It is important to mention that $\dcat{C}(A, \cat{M})$ has 
more structure than just a DG category. Here the objects have DG structures 
too, and there is the cohomology functor 
$\opn{H} : \dcat{C}_{\mrm{str}}(A, \cat{M}) \to 
\dcat{G}_{\mrm{str}}(\cat{M})$,
where the latter is the category of graded objects in $\cat{M}$. 
The cohomology functor determines the {\em quasi-isomorphisms}: these are the 
morphisms $\psi$ in $\dcat{C}_{\mrm{str}}(A, \cat{M})$ s.t.\ 
$\opn{H}(\psi)$ is an isomorphism. 
The set of quasi-isomorphisms is denoted by 
$\dcat{S}(A, \cat{M})$.  

\medskip \noindent 
{\bf Section \ref{sec:stand-tri}.} 
We talk about the {\em translation functor} and {\em standard triangles}
in $\dcat{C}(A, \cat{M})$.
This section consists mostly of new material, some of it implicit in the paper 
\cite{BoKa} on {\em pretriangulated DG categories}. 

The translation $\opn{T}(M)$ of a DG module $M$ is the usual one (a shift by 
$1$ in degree, and differential $-\d_M$). A calculation shows that 
$\opn{T} : \dcat{C}(A, \cat{M}) \to \dcat{C}(A, \cat{M})$
is a DG functor. We introduce the ``little t operator'', which is an invertible 
degree $-1$ morphism 
$\opn{t} : \opn{Id} \to \opn{T}$
of DG functors from $\dcat{C}(A, \cat{M})$ to itself.
The operator $\opn{t}$ facilitates many calculations.  

A morphism $\phi : M \to N$ in $\dcat{C}_{\mrm{str}}(A, \cat{M})$
gives rise to the standard triangle
\[ M \xar{\phi} N \xar{ e_{\phi}} \opn{Cone}(\phi) \xar{p_{\phi}} \opn{T}(M) \]
in $\dcat{C}_{\mrm{str}}(A, \cat{M})$.
As a graded module the standard cone is 
\[ \opn{Cone}(\phi) := N \oplus \opn{T}(M) =
\bmat{N \\ \opn{T}(M)}  , \]
written as a column module. The differential of $\opn{Cone}(\phi)$ is left 
multiplication by the matrix of degree $1$ operators 
\[ \d_{\mrm{Cone}} :=
\bmat{\d_N & \phi \circ \opn{t}_M^{-1} \\[0.1em] 0 & \d_{\opn{T}(M)}} \]

Consider a DG functor
\begin{equation} \label{eqn:1310}
F : \dcat{C}(A, \cat{M}) \to \dcat{C}(B, \cat{N}) , 
\end{equation}
where $B$ is another DG ring and  $\cat{N}$ is another abelian 
category. In Theorem \ref{thm:1150} we show that there is a canonical 
isomorphism
\begin{equation} \label{eqn:1720}
\tau_F : F \circ \opn{T} \iso \opn{T} \circ \, F
\end{equation}
of functors 
$\dcat{C}_{\mrm{str}}(A, \cat{M}) \to \dcat{C}_{\mrm{str}}(B, \cat{N})$,
called the {\em translation isomorphism}. The pair 
$(F, \tau_F)$ is called a {\em T-additive functor}. 
Then, in Theorem \ref{thm:1185}, 
we prove that $F$ sends standard triangles in  
$\dcat{C}_{\mrm{str}}(A, \cat{M})$
to standard triangles in $\dcat{C}_{\mrm{str}}(B, \cat{N})$. 

We end this section with several examples of DG functors.
These examples are prototypes -- they can be easily extended to other setups. 
 

 \medskip \noindent 
{\bf Section \ref{sec:triangulated}.}
We start with the theory of {\em triangulated categories} and 
{\em triangulated functors}, mainly following \cite{RD}. Because the  
octahedral axiom plays no role in our book, we give it minimal attention. 

The {\em homotopy category} $\dcat{K}(A, \cat{M})$ has the same objects as 
$\dcat{C}(A, \cat{M})$, and its morphisms are 
\[ \opn{Hom}_{\dcat{K}(A, \cat{M})} (M, N) := 
\opn{H}^0 \bigl( \opn{Hom}_{\dcat{C}(A, \cat{M})} (M, N) \bigr) . \]
There is a functor 
$\opn{P} : \dcat{C}_{\mrm{str}}(A, \cat{M}) \to \dcat{K}(A, \cat{M})$, 
which is the identity on objects and surjective on morphisms. 

In Subsection \ref{subsec:K-is-triang} we prove that the homotopy 
category $\dcat{K}(A, \cat{M})$ is triangulated. Its {\em distinguished 
triangles} are the triangles that are isomorphic to the images under the 
functor $\opn{P}$ of the standard triangles in 
$\dcat{C}_{\mrm{str}}(A, \cat{M})$. 

Theorem \ref{thm:1265} says that given a DG 
functor $F$ as in (\ref{eqn:1310}), with translation isomorphism $\tau_F$
from (\ref{eqn:1720}), the T-additive functor 
\begin{equation} \label{eqn:5012}
(F, \tau_F) : \dcat{K}(A, \cat{M}) \to \dcat{K}(B, \cat{N})
\end{equation}
is triangulated. This seems to be a new result, unifying well-known yet 
disparate examples. 

In Subsection \ref{subsec:opp-hom-triang} we put a triangulated structure on 
the opposite homotopy category 
$\dcat{K}(A, \cat{M})^{\mrm{op}}$.
We prove that a contravariant DG functor $F$, like (\ref{eqn:1310}) but
with source $\dcat{C}(A, \cat{M})^{\mrm{op}}$, induces a 
contravariant triangulated functor $(F, \tau_F)$, like (\ref{eqn:5012}) but
with source $\dcat{K}(A, \cat{M})^{\mrm{op}}$.



\medskip \noindent 
{\bf Section \ref{sec:loc-cats}.}
In this section we take a close look at {\em localization of categories}. 
We give a detailed proof of the theorem on {\em Ore localization} (also known 
as noncommutative localization). We then prove that the localization 
$\cat{K}_{\cat{S}}$ of a linear category $\cat{K}$ at a denominator set 
$\cat{S}$ is a linear category too, and the localization functor 
$\opn{Q} : \cat{K} \to \cat{K}_{\cat{S}}$ is linear. 


\medskip \noindent 
{\bf Section \ref{sec:der-cat}.}
We begin by proving that if $\cat{K}$ is a triangulated category and 
$\cat{S} \sub \cat{K}$ is a {\em denominator set of cohomological origin}, then 
the localized category $\cat{K}_{\cat{S}}$ is triangulated, 
and the localization functor 
$\opn{Q} : \cat{K} \to \cat{K}_{\cat{S}}$ is triangulated.

In the case of the triangulated category $\dcat{K}(A, \cat{M})$, and the 
quasi-iso\-morph\-isms $\dcat{S}(A, \cat{M})$ in it, we get the {\em derived 
category}
\begin{equation} \label{eqn:5013}
\dcat{D}(A, \cat{M}) := \dcat{K}(A, \cat{M})_{\dcat{S}(A, \cat{M})} ,
\end{equation}
and the triangulated localization functor 
\begin{equation} \label{eqn:5014}
\opn{Q} : \dcat{K}(A, \cat{M}) \to \dcat{D}(A, \cat{M}) . 
\end{equation}

We look at the full subcategory $\dcat{K}^{\star}(A, \cat{M})$
of $\dcat{K}(A, \cat{M})$ corresponding to a boundedness condition 
$\star \in \{ +, -, \mrm{b} \}$. 
We prove that when the DG ring $A$ is nonpositive, the localization
$\dcat{D}^{\star}(A, \cat{M})$ of $\dcat{K}^{\star}(A, \cat{M})$
with respect to the quasi-iso\-morphisms in it embeds
fully faithfully in $\dcat{D}(A, \cat{M})$. 
We also prove that the obvious functor 
$\cat{M} \to \dcat{D}(\cat{M})$ is fully faithful.

The section ends with a study of the triangulated structure of the opposite 
derived category $\dcat{D}(A, \cat{M})^{\mrm{op}}$.

\medskip \noindent 
{\bf Section \ref{sec:der-funcs}.}
In this section we talk about {\em derived functors}. To make the definitions 
of derived functors precise, we introduce some {\em $2$-categorical 
notation} here (in Subsections \ref{subsec:2-cat-not}-\ref{subsec:fun-cats}).

In Subsection \ref{subsec:abstr-der-funs} we look at {\em abstract derived 
functors}. Namely, $\cat{K}$ and $\cat{E}$ are categories (without any extra 
structure), $F : \cat{K} \to \cat{E}$ is a functor, and 
$\cat{S} \sub \cat{K}$ is a multiplicatively closed set of morphisms. 
The localization of $\cat{K}$ w.r.t.\ $\cat{S}$ is 
$\opn{Q} : \cat{K} \to \cat{K}_{\cat{S}}$. 
A {\em right derived functor of $F$ w.r.t.\ $\cat{S}$} is a pair 
$(\mrm{R} F ,\eta^{\mrm{R}})$, where 
$\mrm{R} F : \cat{K}_{\cat{S}} \to \cat{E}$
is a functor, and 
$\eta^{\mrm{R}} : F \to \mrm{R} F \circ \opn{Q}$
is a morphism of triangulated functors. The pair $(\mrm{R} F ,\eta^{\mrm{R}})$ 
has a universal property (it is initial among all such pairs), making it unique 
up to a unique isomorphism. The  {\em left derived functor}
$(\mrm{L} F ,\eta^{\mrm{L}})$ is defined similarly.

We provide a general existence theorem for derived functors. For the right 
derived functor $\mrm{R} F$, existence is proved when $\cat{S}$ is a left 
denominator set, and there exists a full 
subcategory $\cat{J} \sub \cat{K}$ which is $F$-acyclic and 
right-resolves all objects of $\cat{K}$.  Likewise, for the left derived functor
$\mrm{L} F$ we prove existence when $\cat{S}$ is a right denominator set, and 
there exists a full subcategory $\cat{P} \sub \cat{K}$ which is $F$-acyclic and 
left-resolves all objects of $\cat{K}$. 

In Subsection \ref{subsec:tri-der-funs} we specialize to {\em triangulated 
derived functors}. Here $\cat{K}$ and $\cat{E}$ are triangulated categories, 
$F : \cat{K} \to \cat{E}$ is a triangulated functor, and 
$\cat{S} \sub \cat{K}$ is a multiplicatively closed set of cohomological 
origin (i.e.\ $\cat{S}$ consists of the quasi-isomorphisms w.r.t.\ 
some cohomological functor $H : \cat{K} \to \cat{M}$). The right and left 
derived functors 
$\mrm{R} F, \mrm{L} F : \cat{K}_{\cat{S}} \to \cat{E}$
are defined like in the abstract setting, and their uniqueness 
is also proved the same way. Existence requires resolving subcategories 
$\cat{J}, \cat{P} \sub \cat{K}$ as above, that are also full triangulated 
subcategories. 

The section is concluded with a discussion of {\em contravariant triangulated 
derived functors}. 

\medskip \noindent 
{\bf Section \ref{sec:DG-tri-bifun}.}
This section is devoted to  {\em DG and triangulated bifunctors}. 
We prove that a DG bifunctor 
\[ F : \dcat{C}(A_1, \cat{M}_1) \times \dcat{C}(A_2, \cat{M}_2)
\to \dcat{C}(B, \cat{N}) \]
induces a triangulated bifunctor 
\[ (F, \tau_1, \tau_2) :  \dcat{K}(A_1, \cat{M}_1) \times 
\dcat{K}(A_2, \cat{M}_2) \to \dcat{K}(B, \cat{N}) . \]

Then we define left and right derived bifunctors, and prove their uniqueness 
and existence, under suitable conditions. Bifunctors that are contravariant in 
one or two of the arguments are also studied. 


\medskip \noindent 
{\bf Section \ref{sec:resol}.}
Here we define {\em K-injective} and {\em K-projective} objects in 
$\dcat{K}(A, \cat{M})$, and also  {\em K-flat} objects in 
$\dcat{K}(A)$. These constitute full triangulated subcategories of
$\dcat{K}(A, \cat{M})$, and we refer to them as {\em resolving subcategories}.
The category $\dcat{K}^{\star}(A, \cat{M})_{\mrm{inj}}$ of K-injectives in 
$\dcat{K}^{\star}(A, \cat{M})$, for a boundedness condition $\star$, 
plays the role of the category $\cat{J}$ in Section \ref{sec:der-funcs} 
above, and the category of K-projectives 
$\dcat{K}^{\star}(A, \cat{M})_{\mrm{prj}}$ 
plays the role of the category $\cat{P}$ there.
The K-flat DG modules are acyclic for the tensor functor. 

Furthermore, we prove that the functors 
$\opn{Q} : \dcat{K}^{\star}(A, \cat{M})_{\mrm{inj}} \to 
\dcat{D}^{\star}(A, \cat{M})$
and 
$\opn{Q} : \dcat{K}^{\star}(A, \cat{M})_{\mrm{prj}} \to 
\dcat{D}^{\star}(A, \cat{M})$,
see equation (\ref{eqn:5014}), are fully faithful. 


\medskip \noindent 
{\bf Section \ref{sec:exist-resol}.}
In this section we prove {\em existence} of K-injective, K-projective and 
K-flat resolutions in several important cases of 
$\dcat{C}^{\star}(A, \cat{M})$~:
\begin{itemize}
\item K-projective resolutions in $\dcat{C}^{-}(\cat{M})$, where $\cat{M}$ is 
an abelian category with enough projectives. This is classical (i.e.\ it is 
already in \cite{RD}).  

\item K-projective resolutions in $\dcat{C}^{}(A)$, where $A$ is any DG ring. 
This includes the category of unbounded complexes of modules over a ring $A$. 

\item K-injective resolutions in $\dcat{C}^{+}(\cat{M})$, where $\cat{M}$ is 
an abelian category with enough injectives. This is classical too. 

\item K-injective resolutions in $\dcat{C}^{}(A)$, where $A$ is any DG ring. 
\end{itemize}
Our proofs are explicit, and we use limits of complexes cautiously (since 
this is known to be a pitfall). We follow several sources: \cite{RD}, 
\cite{Spa}, \cite{Kel1}, \cite{KaSc1} and private notes provided by B. Keller.


\medskip \noindent 
{\bf Section \ref{sec:adj-equ-cohdim}.}
This section is quite varied. 
In Subsections \ref{subsec:RHom}-\ref{subsec:LTors} there is
a detailed look at the derived bifunctors 
$\opn{RHom}(-,-)$ and $(- \ot^{\mrm{L}} -)$. 

Next, in Subsection \ref{subsec:way-out}, we study {\em cohomological 
dimensions of functors}. This is a refinement of the notion of {\em way-out 
functor} from \cite{RD}. It is used in Subsection \ref{subsec:thm-funcs-fin} to 
prove a few theorems about triangulated functors, such as a sufficient condition 
for a morphism $\ze : F \to G$ of triangulated functors to be an isomorphism. 

In Subsections \ref{subsec:adjs-NC}-\ref{subsec:DGR-quisoms} we study 
several adjunction formulas that involve 
the bifunctors $\opn{RHom}(-,-)$ and $(- \ot^{\mrm{L}} -)$.
We define {\em derived forward adjunction} and  
{\em derived backward adjunction}. These adjunction formulas hold for 
arbitrary DG rings (without any commutativity or boundedness conditions). 
We prove that if $A \to B$ is a {\em quasi-isomorphism of DG rings}, then 
the restriction functor 
$\opn{Rest}_{B / A} : \dcat{D}(B) \to \dcat{D}(A)$
is an equivalence, and it respects the derived bifunctors 
$\opn{RHom}(-,-)$ and $(- \ot^{\mrm{L}} -)$.

Resolutions of DG rings are important in several contexts. In Subsection 
\ref{subsec:DG-ring-res} we provide a full proof that given a 
DG $\K$-ring $A$, there exists a {\em noncommutative semi-free DG ring 
resolution} $\til{A} \to A$ over $\K$.

In Subsection \ref{subsec:der-tens-eval} there is a very general 
theorem providing sufficient conditions for the {\em derived tensor-evaluation 
morphism}
\begin{equation} \label{eqn:4610}
\opn{ev}^{\mrm{R, L}}_{L, M, N} :
\opn{RHom}_{A}(L, M) \ot^{\mrm{L}}_{B} N \to 
\opn{RHom}_{A}(L, M \ot^{\mrm{L}}_{B} N)
\end{equation}
in $\dcat{D}(\K)$ to be an isomorphism. Here $A$ and $B$ are DG $\K$-rings, 
$L \in \dcat{D}(A)$, $M \in \dcat{D}(A \ot_{\K} B^{\mrm{op}})$ and 
$N \in \dcat{D}(B)$.

Lastly, in Subsection \ref{subsec:adjs-comm} we present some 
adjunction formulas that pertain only to weakly commutative DG rings. 




\medskip \noindent 
{\bf Section \ref{sec:dual-cplx-comm-rng}.}
This section starts (in the first three subsections) by retelling the material 
in \cite{RD} on {\em dualizing complexes} and {\em residue complexes} over 
noetherian commutative rings.

A complex $R \in \dcat{D}(A)$ is called {\em dualizing} if it has finitely 
generated cohomology modules, finite injective dimension, and the {\em derived 
Morita property}, which says that the {\em derived homothety morphism}
\begin{equation} \label{eqn:4670}
\opn{hm}^{\mrm{R}}_{R} : A \to \opn{RHom}_{A}(R, R)
\end{equation}
in $\dcat{D}(A)$ is an isomorphism. 
(This is a variant of the derived Hom-evaluation morphism 
$\opn{ev}^{\mrm{R}}_M$ from equation (\ref{eqn:3248}), with 
$M = A$.)

To a dualizing complex $R$ we associate the {\em duality functor}
\[ D := \opn{RHom}_{A}(-, R) : \dcat{D}(A)^{\mrm{op}}
\to \dcat{D}(A) . \]
The duality functor induces an equivalence of triangulated categories
$D : \dcat{D}_{\mrm{f}}(A)^{\mrm{op}} \lb \to \dcat{D}_{\mrm{f}}(A)$.
Here $\dcat{D}_{\mrm{f}}(A)$ is the full subcategory of $\dcat{D}(A)$
on the complexes with finitely generated cohomology modules. 

A residue complex is a dualizing complex 
$\KK$ that consists of injective modules of the correct multiplicity in each 
degree. If $R = \KK$ is a residue complex, then the associated duality functor 
is $D = \opn{Hom}_{A}(-, \KK)$.

We prove uniqueness of dualizing complexes over a noetherian commutative ring 
$A$ (up to the obvious twists), and existence when $A$ is {\em essentially 
finite type over a regular noetherian ring} $\K$. In this book the adjective 
``regular noetherian'' includes the condition that $\K$ has finite Krull 
dimension. 

There is a stronger uniqueness property for residue complexes. 
Residue complexes exist whenever dualizing complexes exist: given a dualizing 
complex $R$, its minimal injective resolution $\KK$ is a residue complex. 
To understand residue complexes we review the {\em Matlis classification of 
injective $A$-modules}.  

In remarks we provide sketches of Matlis Duality, Local Duality and the 
interpretation of Cohen-Macaulay complexes as perverse modules.  

In the last two subsections we talk about {\em Van den Bergh rigidity}.
Let $\K$ be a regular noetherian ring, and let $A$ be a flat essentially finite 
type (FEFT) $\K$-ring. (The flatness condition is just to simplify matters; see 
Remark  \ref{rem:2285}.)  Given a complex $M \in \dcat{D}(A)$, its {\em square 
relative to $\K$} is the complex
\[ \opn{Sq}_{A / \K}(M) := 
\opn{RHom}_{A \ot_{\K} A}(A, M \ot^{\mrm{L}}_{\K} M) \in \dcat{D}(A) . \]
We prove that $\opn{Sq}_{A / \K}$ is a {\em quadratic functor}. 
A {\em rigid complex over $A$ relative to $\K$} is a pair 
$(M, \rho)$, where $M \in \dcat{D}(A)$ and 
$\rho : M \iso \opn{Sq}_{A / \K}(M)$
is an isomorphism in $\dcat{D}(A)$, called a {\em rigidifying isomorphism}.
Given another rigid complex $(N, \si)$, a {\em rigid morphism} between 
them is a morphism $\phi : M \to N$ in $\dcat{D}(A)$ such that 
the diagram 
\begin{equation} \label{eqn:5015}
\UseTips \xymatrix @C=8ex @R=6ex {
M
\ar[d]_{\phi}
\ar[r]^(0.4){\rho}
&
\opn{Sq}_{A / \K}(M)
\ar[d]^{\opn{Sq}_{A / \K}(\phi)}
\\
N
\ar[r]^(0.4){\si}
&
\opn{Sq}_{A / \K}(N)
}
\end{equation}
in $\dcat{D}(A)$ is commutative. 
We prove that if $(M, \rho)$ is a rigid complex such that $M$ has the derived 
Morita property, then the only rigid automorphism of $(M, \rho)$ is the 
identity. 

A {\em rigid dualizing complex over $A$ relative to $\K$} is a rigid complex 
$(R, \rho)$ such that $R$ is a dualizing complex. We prove that if $A$ has a 
rigid dualizing complex, then it is unique up to a unique rigid isomorphism.
Existence of a rigid dualizing complex is harder to prove, and we just give a 
reference to it. 

A {\em rigid residue complex over $A$ relative to $\K$} is a rigid complex 
$(\KK, \rho)$ such that $\KK$ is a residue complex. These always exist, 
and they are unique in the following very strong sense: if $(\KK', \rho')$ 
is another rigid residue complex, then there is a 
unique isomorphism $\phi : \KK \iso \KK'$ in $\dcat{C}_{\mrm{str}}(A)$, such 
that $\opn{Q}(\phi)$ is a rigid isomorphism in $\dcat{D}(A)$.

We end this section 
with two remarks (Remark \ref{rem:4192} and \ref{rem:4193}) 
that explain how rigid residue complexes allow a new approach to {\em residues 
and duality on schemes and Deligne-Mumford stacks}, with references. 


\medskip \noindent 
{\bf Section \ref{sec:perf-tilt-NC}.}
We begin (in Subsection \ref{subsec:perf-dgmods}) with a systematic 
study of {\em algebraically perfect DG modules} over a NC DG ring $A$. 
The abbreviation ``NC'' stands for ``noncommutative''. By 
definition, a DG module $L \in \dcat{D}(A)$ is called algebraically perfect if 
it belongs to the saturated full triangulated subcategory of $\dcat{D}(A)$ 
generated by the DG module $A$. We give several characterizations of 
algebraically perfect DG modules; one of them is that $L$ is a {\em compact 
object} of $\dcat{D}(A)$. When $A$ is a ring, we prove that $L$ is  
algebraically perfect if and only if it is isomorphic, in $\dcat{D}(A)$, to a 
bounded complex of finitely generated projective $A$-modules. 

In Subsection \ref{subsec:der-morita} we prove the following {\em 
Derived Morita Theorem}, Theorem \ref{thm:3430}:
Let $\cat{E} \sub \dcat{D}(A)$ be a full triangulated subcategory which is 
closed under infinite direct sums, let $P$ be a compact generator of 
$\cat{E}$ which is either K-projective or K-injective, and let
$B := \opn{End}_{A}(P)^{\mrm{op}}$. Then the functor 
$\opn{RHom}_{A}(P, -) : \cat{E} \to \dcat{D}(B)$
is an equivalence of triangulated categories, with quasi-inverse 
$P \ot^{\mrm{L}}_{B} (-)$. 

From Subsection \ref{subsec:flat-dg-rng} to the end of this section 
we assume that the DG rings in question are K-flat over the base ring $\K$. 
Note that every DG $\K$-ring $A$ admits a NC semi-free DG ring resolution 
$\til{A} \to A$. The DG ring $\til{A}$ is K-flat over $\K$, and the 
restriction functor $\dcat{D}(A) \to \dcat{D}(\til{A})$ is an equivalence of 
triangulated categories; so the flatness restriction can be  
circumvented. Subsection \ref{subsec:flat-dg-rng} contains some basic 
constructions of derived functors between categories of DG bimodules (that 
require the flatness over the base ring). 

Next, in Subsection \ref{subsec:tilting}, we define {\em tilting DG 
bimodules}. A DG bimodule $T \in \dcat{D}(B \ot_{\K} A^{\mrm{op}})$ is called
a tilting DG $B$-$A$-bimodule if there exists some 
$S \in \lb \dcat{D}(A \ot_{\K} B^{\mrm{op}})$ 
with isomorphisms $S \ot^{\mrm{L}}_{B} T \cong A$ in 
$\dcat{D}(A^{\mrm{en}})$ and $T \ot^{\mrm{L}}_{A} S \cong B$ in 
$\dcat{D}(B^{\mrm{en}})$. Here 
$A^{\mrm{en}} := A \ot_{\K} A^{\mrm{op}}$ is the {\em enveloping DG ring of 
$A$}, and likewise $B^{\mrm{en}}$. Among other results, we prove that 
$T \in \dcat{D}(B \ot_{\K} A^{\mrm{op}})$
is tilting if and only if $T$ is a compact generator on the $B$ side
(i.e.\ it is a compact generator of $ \dcat{D}(B)$), and it 
has the {\em NC derived Morita property} on the $B$ side, namely the canonical
morphism $A \to \opn{RHom}_{B}(T, T)$
in $\dcat{D}(A^{\mrm{en}})$ is an isomorphism.
We also prove the {\em Rickard-Keller Theorem}, asserting that if $A$ and $B$ 
are rings, and there exists a $\K$-linear equivalence of 
triangulated categories
$\dcat{D}(A) \to \dcat{D}(B)$, then there exists a tilting DG $B$-$A$-bimodule. 

Lastly, in Subsection \ref{subsec:tilting-rings}, we introduce the {\em NC
derived Picard group} $\opn{DPic}_{\K}(A)$ of a flat $\K$-ring $A$. The 
structure of this group is calculated when $A$ is either local or commutative.  


\medskip \noindent 
{\bf Section \ref{sec:alg-gra-rings}.}
This section, as well as Sections \ref{sec:der-tors-conn} and \ref{sec:BDC}, 
are on {\em algebraically graded rings}, which is our name for 
$\Z$-graded rings that have lower indices and do not involve the Koszul sign 
rule.  (This is in contrast with the cohomologically graded rings mentioned 
above, that underly DG rings). Simply put, these are the usual graded rings 
that one encounters in textbooks on commutative and noncommutative algebra. 
With few exceptions, the base ring $\K$ in the four final sections of the book 
is a field. 

Let $A = \boplus_{i \in \Z} A_i$ be an algebraically graded central $\K$-ring. 
In Section \ref{sec:alg-gra-rings} we define the category 
of algebraically graded $A$-modules $\dcat{M}(A, \mrm{gr})$. 
Its objects are the algebraically graded (left) $A$-modules
$M = \boplus_{i \in \Z} M_i$, and the homomorphisms are the degree $0$
$A$-linear homomorphisms. It is a $\K$-linear abelian category. 
We talk about finiteness in the algebraically graded context, and about 
various kinds of homological properties, such as {\em 
graded-injectivity}. 

The category of complexes with entries in $\dcat{M}(A, \mrm{gr})$ is the DG 
category \lb 
$\dcat{C}_{}(A, \mrm{gr}) := 
\dcat{C}_{} \bigl( \dcat{M}(A, \mrm{gr}) \bigr)$.
Its objects are bigraded: a complex 
$M \in \dcat{C}_{}(A, \mrm{gr})$ has a direct sum decomposition 
$M = \boplus_{i, j} M^i_j$ into $\K$-modules. Here $i$ is the cohomological 
degree and $j$ is the algebraic degree. The differential goes like this:
$\d_M : M^i_j \to M^{i + 1}_j$. 

Just like for any other abelian category, we have the derived category 
$\dcat{D}_{}(A, \mrm{gr}) := 
\dcat{D}_{} \bigl( \dcat{M}(A, \mrm{gr}) \bigr)$.
This is a triangulated category. We present (quickly) the algebraically graded 
variants of K-projective resolutions etc., and the relevant derived functors.

Special emphasis is given to {\em connected graded rings}. An algebraically 
graded ring $A$ is called connected if $A = \boplus_{i \in \N} A_i$, 
$A_0 = \K$, and each $A_i$ is a finite $\K$-module.
In a connected graded ring $A$ there is the {\em augmentation ideal} 
$\m := \boplus_{i > 0} A_i$. We view $A / \m \cong \K$ as a graded 
$A$-bimodule. 

Among the connected graded rings we single out the {\em Artin-Schelter regular 
graded rings}. A noetherian connected graded ring $A$ is called AS regular if 
it has finite graded global dimension, and if there are isomorphisms
\begin{equation} \label{eqn:4621}
\opn{RHom}_{A}(\K, A) \cong \opn{RHom}_{A^{\mrm{op}}}(\K, A) \cong
\K(l)[-n]
\end{equation}
in $\dcat{D}(\K)$ for some integers $l$ and $n$. One of the results is that if 
$A$ is a noetherian connected graded ring, $a \in A$ is a regular central 
homogeneous element of positive degree, and the ring $B := A / (a)$ is AS 
regular, then $A$ is also AS regular. 


\medskip \noindent 
{\bf Section \ref{sec:der-tors-conn}.}
Let $A$ be a connected graded ring over the base field $\K$, with augmentation 
ideal $\m$. In this section we study {\em derived 
$\m$-torsion}, both for complexes of graded $A$-modules and for complexes of 
graded bimodules. By this we mean that, taking a second graded ring $B$, we 
look at the triangulated functor 
\[ \mrm{R} \Ga_{\m} : \dcat{D}(A \ot_{\K} B^{\mrm{op}}, \mrm{gr}) \to
\dcat{D}(A \ot_{\K} B^{\mrm{op}}, \mrm{gr}) .  \]

One of the main results (Theorem \ref{thm:3750}) says that if $A$ is 
a left noetherian connected graded ring, and if the functor 
$\mrm{R} \Ga_{\m}$ has finite cohomological dimension, then there is a 
functorial isomorphism 
\begin{equation} \label{eqn:5016}
\opn{ev}^{\mrm{R, L}}_{\m, M} : 
P \ot^{\mrm{L}}_{A} M \iso \mrm{R} \Ga_{\m}(M)  
\end{equation}
for $M \in \dcat{D}(A \ot_{\K} B^{\mrm{op}}, \mrm{gr})$,
where
\begin{equation} \label{eqn:4620}
P := \mrm{R} \Ga_{\m}(A) \in \dcat{D}(A^{\mrm{en}}, \mrm{gr}) . 
\end{equation}
We also prove the {\em NC MGM Equivalence} in the connected graded context (see 
Theorem \ref{thm:3792}).

The {\em $\chi$ condition} of M. Artin and J.J. Zhang is introduced in 
Subsection \ref{subsec:sym-der-tor}. We study how this condition 
interacts with derived torsion.
To a complex of graded $A$-bimodules, namely an object of 
$\dcat{D}(A^{\mrm{en}}, \mrm{gr})$, we can also apply derived $\m$-torsion from 
the right side, thus obtaining a complex 
$\mrm{R} \Ga_{\m^{\mrm{op}}}(M) \in \dcat{D}(A^{\mrm{en}}, \mrm{gr})$. 
Theorem \ref{thm:4051} says that if $A$ is a 
noetherian connected graded ring of finite local cohomological dimension (i.e.\ 
both functors $\mrm{R} \Ga_{\m}$ and $\mrm{R} \Ga_{\m^{\mrm{op}}}$ have finite 
cohomological dimension), that satisfies the $\chi$ condition, then there is a 
functorial isomorphism 
\begin{equation} \label{eqn:5017}
\ep_M : \mrm{R} \Ga_{\m}(M) \iso \mrm{R} \Ga_{\m^{\mrm{op}}}(M) 
\end{equation}
in $\dcat{D}(A^{\mrm{en}}, \mrm{gr})$ for all complexes 
$M \in \dcat{D}(A^{\mrm{en}}, \mrm{gr})$ whose cohomologies 
are finite $A$-modules on both sides. We call this phenomenon {\em symmetric 
derived $\m$-torsion}. 


\medskip \noindent 
{\bf Section \ref{sec:BDC}.}
The focus of this section is on {\em balanced NC dualizing complexes}, 
following \cite{Ye1}.
Let $A$ be a noetherian connected graded ring over the base field $\K$. A 
complex $R \in \dcat{D}^{\mrm{b}}(A^{\mrm{en}}, \mrm{gr})$
is called a {\em graded NC dualizing complex} if it satisfies these three 
conditions:
\begin{enumerate}
\rmitem{i} The bimodules $\opn{H}^q(R)$ are finitely generated 
$A$-modules on both sides.

\rmitem{ii} The complex $R$ has finite graded injective 
dimension on both sides.

\rmitem{iii} The complex $R$ has the NC derived Morita property on both 
sides; namely the canonical morphisms 
$A \to \opn{RHom}_{A}(R, R)$ and 
$A \to \opn{RHom}_{A^{\mrm{op}}}(R, R)$ 
in $\dcat{D}(A^{\mrm{en}}, \mrm{gr})$ are isomorphisms. 
\end{enumerate}

A balanced NC dualizing complex over $A$ is a pair $(R, \be)$,
where $R$ is a graded NC dualizing complex over $A$ with symmetric derived 
$\m$-torsion, and 
$\be : \mrm{R} \Ga_{\m}(R) \iso A^*$
is an isomorphism in $\dcat{D}(A^{\mrm{en}}, \mrm{gr})$, called a {\em 
balancing isomorphism}. Here and later we write 
$M^* := \opn{Hom}_{\K}(M, \K)$, the graded $\K$-linear dual of a graded module 
$M$. The graded bimodule $A^*$ is a graded-injective $A$-module on both sides. 

To a balanced dualizing complex $R$ we associate the duality functors 
\[ D_A := \opn{RHom}_{A}(-, R) : \dcat{D}(A, \mrm{gr})^{\mrm{op}}
\to \dcat{D}(A^{\mrm{op}}, \mrm{gr}) \]
and 
\[ D_{A^{\mrm{op}}} := \opn{RHom}_{A^{\mrm{op}}}(-, R) : 
\dcat{D}(A^{\mrm{op}}, \mrm{gr})^{\mrm{op}}
\to \dcat{D}(A, \mrm{gr}) . \]
When restricted to the subcategories of complexes with finitely generated 
cohomology modules, this gives an equivalence 
\begin{equation} \label{eqn:5018}
D_A :  \dcat{D}_{\mrm{f}}(A, \mrm{gr})^{\mrm{op}}
\to \dcat{D}_{\mrm{f}}(A^{\mrm{op}}, \mrm{gr})
\end{equation}
with quasi-inverse $D_{A^{\mrm{op}}}$. 

A balanced dualizing complex $(R, \be)$ is unique up to a unique isomorphism 
(Theorem \ref{thm:3710}), and it satisfies the NC Local Duality Theorem 
\ref{thm:4616}.  Results of Yekutieli, Zhang and M. Van den Bergh say that 
a noetherian connected graded ring $A$ has a balanced dualizing complex if and 
only if $A$ satisfies the $\chi$ condition and has finite local cohomological 
dimension. The formula for the balanced dualizing complex is 
$R := P^*$, where $P$ is the complex from (\ref{eqn:4620}). 
See Corollary \ref{cor:4585}.

If $A$ is an AS regular (or more generally AS Gorenstein) graded ring, then it 
has a balanced dualizing complex $R = A(\phi, -l)[n]$. Here $l$ and $n$ are the 
integers from formula (\ref{eqn:4621}), and $\phi$ is a graded ring 
automorphism 
of $A$. See Corollary \ref{cor:3780}. 



\medskip \noindent 
{\bf Section \ref{sec:rigid-DC-NC}.}
The final section of the book deals with {\em NC rigid dualizing complexes}. 
Let $A$ be a NC central $\K$-ring (where $\K$ is the base field). A {\em NC 
dualizing complex} over $A$ is a complex 
$R \in \dcat{D}^{\mrm{b}}(A^{\mrm{en}})$
whose cohomology bimodules are finitely generated on both sides; it has finite 
injective dimension on both sides; and it has the NC derived Morita property on 
both sides. This is a modification of the graded definition given above.

Let $M \in \dcat{D}(A^{\mrm{en}})$. The {\em NC square} of $M$
is the complex 
\begin{equation} \label{eqn:4660}
\opn{Sq}_{A / \K}(M) := \opn{RHom}_{A^{\mrm{en}}}(A, R \ot_{\K} R) \in 
\dcat{D}(A^{\mrm{en}}) .
\end{equation}
This formula is ambiguous, because  
the complex of $\K$-modules $R \ot_{\K} R$ has on it four distinct 
commuting actions by the ring $A$; but formula (\ref{eqn:4660}) is made precise 
at the beginning of Subsection \ref{subsec:RNCDC-uniq}. 
A {\em rigidifying isomorphism} for $M$ is an isomorphism 
$\rho : M \iso  \opn{Sq}_{A / \K}(M)$
in $\dcat{D}(A^{\mrm{en}})$, and the pair $(M, \rho)$ is a {\em NC rigid 
complex}. A NC rigid dualizing complex is a NC rigid complex $(R, \rho)$ 
such that $R$ is a NC dualizing complex. 

The definition of NC rigid dualizing complex was introduced by Van den Bergh in 
his paper \cite{VdB}. He also proved that a rigid NC dualizing complex is 
unique up to isomorphism. Later, in \cite{Ye4}, it was proved that this 
isomorphism is itself unique. We reproduce these results in Subsection 
\ref{subsec:RNCDC-uniq}.  

We also give Van den Bergh's theorem on the existence of 
NC rigid dualizing complexes from \cite{VdB}, with a full proof. 
This is Theorem \ref{thm:4396}. Here is what it says: 
Suppose the ring $A$ admits a filtration 
$F = \{ F_j(A) \}_{j \geq -1}$
such that the associated graded ring $\opn{Gr}^F(A)$ 
is noetherian connected and has a balanced dualizing complex. Then $A$ has a 
rigid dualizing complex. 

An important special case of the Van den Bergh Existence Theorem is  
Theorem \ref{thm:4490}. It says that if the graded ring  $\opn{Gr}^F(A)$ from 
the previous paragraph is AS regular, then the rigid NC dualizing complex of 
$A$ is $R = A(\mu)[n]$,
where $\mu$ is a $\K$-ring automorphism of $A$ that respects the filtration 
$F$, and $n$ is an integer. The automorphism $\nu := \mu^{-1}$ is called the 
{\em Nakayama automorphism} of $A$. In modern terminology the ring $A$ is 
called an {\em $n$-dimensional twisted Calabi-Yau ring}. 

Finally we state and prove the {\em Van den Bergh Duality Theorem} for 
Hoch\-schild (co)homology, and give an example of a {\em Calabi-Yau category} 
of fractional dimension. 



\mysubsection{What Is Not in the Book}
A very important aspect of the theory of derived categories is {\em geometric}. 
Unfortunately our book does not discuss this aspect (except in passing).
We had hoped to include the geometric aspect, but as the 
book grew longer this became impractical. 

Geometric derived categories come in two kinds. The first kind is the derived 
category $\dcat{D}(X) = \dcat{D}(\cat{Mod} \OO_X)$, 
where $(X, \OO_X)$ is a scheme, and $\cat{Mod} \OO_X$ is the abelian category 
of sheaves of $\OO_X$-modules. This kind of derived category is the subject of 
the original book \cite{RD}. For recent treatments see \cite{Huy}, 
\cite{SP} or \cite{Ols}. The last two references also treat derived 
category $\dcat{D}(\mfrak{X})$ of sheaves of modules on an algebraic stack 
$\mfrak{X}$. We address this aspect of derived categories in Remarks 
\ref{rem:4675}, \ref{rem:4676}, \ref{rem:4192} and \ref{rem:4193}. 

The second kind of geometric derived category is 
$\dcat{D}(\K_X) = \dcat{D}(\cat{Mod} \K_X)$, where $X$ is a topological space
(or a site, such as the \'etale site of a scheme), $\K_X$ is the constant 
sheaf of rings $\K$ on $X$, and 
$\cat{Mod} \K_X$ is the abelian category of sheaves of $\K_X$-modules on $X$.
For this kind of derived categories we recommend the book \cite{KaSc1}. 

One should also mention, in this context, the new and important theories of 
{\em derived algebraic geometry} (DAG), in which schemes are replaced by
{\em derived stacks}. These new geometric objects carry derived categories of 
modules. See the references \cite{Lur}, \cite{To2} and \cite{nLab}, 
Remark \ref{rem:4960} and Example \ref{exa:4205}.
The {\em commutative DG rings} that we discuss briefly in Subsection 
\ref{subsec:DG-rng-mod} are {\em affine derived schemes} from the point of view 
of DAG. 
 

\mysubsection{Prerequisites and Recommended Bibliography}

In preparing the book, the assumptions were that the reader is already familiar 
with these topics:
\begin{itemize}
\item Categories and functors, and classical homological 
algebra, namely the derived functors $\mrm{R}^q F$ and $\mrm{L}_q F$
of an additive functor $F : \cat{M} \to \cat{N}$ between abelian  
categories. For these topics we recommend the books 
\cite{Mac1}, \cite{HiSt} and \cite{Rot}. 

\item Commutative and noncommutative ring theory, e.g.\ from the books
\cite{Eis}, \cite{Mats}, \cite{AlKl}, \cite{Row} and \cite{Rot}.  
\end{itemize}

For the topics above we merely review the material, and point to references 
when necessary.

There are a few earlier books that deal with derived categories, in varying 
degrees of detail and depth. Some of them -- \cite{RD} (the original 
reference), \cite{KaSc1} and \cite{KaSc2} -- served as sources for us when 
writing the present book. Other books, such as \cite{GeMa},  
\cite{We} and \cite{Huy}, are somewhat sketchy in their treatment of derived 
categories, but they might be useful for a reader who wants another perspective 
on the subject. 

Finally we want to mention the evolving online reference \cite{SP}, that 
contains a huge amount of information on all the topics listed above. 

\mysubsection{Credo,  Writing Style and Goals}
Since its inception around 1960, there has been very little literature on the 
theory of derived categories. In some respect, the only detailed account for 
many years was the original book \cite{RD}, written by R. Hartshorne, following 
notes by J.-L. Verdier (for Chapter I) and by A. Grothendieck (for Chapters 
II-VII). Several accounts appeared later as parts of the books 
\cite{We}, \cite{GeMa}, \cite{KaSc1}, \cite{KaSc2}, \cite{Huy}, and maybe a few 
others -- but none of these accounts provided enough detailed content to make 
it possible for a mathematician to learn  how to work with derived categories, 
beyond a rather superficial level. {\em The theory thus remained mysterious}. 

A personal belief of mine is that {\em mathematics should not be mysterious}. 
Some mathematics is very easy to explain. However, a few branches of 
mathematics are truly hard; among them are algebraic geometry and derived 
categories. My moral goal in this book is to demonstrate 
that {\em the theory of derived categories is difficult, but not mysterious}. 

The  series of books \cite{EGA} by Grothendieck and J. Dieudonn\'e, and then 
the book \cite{Har} by Hartshorne, have shown us that {\em algebraic geometry 
is difficult but not mysterious}. The definitions and the statements are 
precise, and the proofs are available (to be read or to be taken on trust, as 
the reader prefers). I hope that the present book will do the same for derived 
categories. (Although I doubt I can match the excellent writing talent of the 
aforementioned authors!).

In more practical terms, the goal of this book is to develop the theory of 
derived categories in a systematic fashion, with full details, and with several 
important applications. The expectation is that our book will open the doors 
for researchers in algebra and geometry (and related disciplines such as 
mathematical physics) to productive work using derived categories. Doors 
that have been to a large extent locked until now (or at least hidden by the 
shrouds of mystery, to use the prior metaphor). 

This book is by far too lengthy for a one-semester graduate course.
(In my rather sluggish lecturing style, Sections 1-13 of the book took about 
four semesters.) The book is intended to be used as a reference, or for 
personal learning. Perhaps a lecturer who has the ability to concentrate the 
material sufficiently, or to choose only certain key aspects, might extract a 
one-semester course from this book; if so, please let me know!



\mysubsection{Acknowledgments}

As already mentioned, the book originated in a course on derived categories, 
that was held at Ben Gurion University in Spring 2012. I want to thank the 
participants of this course for correcting many of my mistakes
(both in real time during the lectures, and afterwards when writing the notes 
\cite{Ye6}).  Thanks also to Joseph Lipman, Pierre Schapira, Amnon Neeman 
and Charles Weibel for helpful discussions during the preparation of that 
course. Vincent Beck, Yang Han and Lucas Simon sent me corrections and useful 
comments on the material in \cite{Ye6}.  

I started writing the book itself while teaching a four
semester course on the subject at Ben Gurion University, spanning the 
academic years 2015-16 and 2016-17. I wish to thank the participants of this 
course, and especially Rishi Vyas, Stephan Snigerov, Asaf Yekutieli, S.D. 
Banerjee, Alberto Boix and William Woods, who contributed material and 
corrected numerous  mistakes. 

Ben Gurion University was generous enough to permit this long project. The 
project was supported by the Israel Science Foundation grant no.\ 253/13. 

Stephan Snigerov, William Woods and Rishi Vyas helped me prepare the manuscript 
for publication, and I wish to thank them for that. Bernhard Keller kindly sent 
me private notes containing clean proofs of several theorems on the 
existence of resolutions of DG modules, and he also helped me with numerous 
suggestions. 

In the process of writing the book I have also benefited from the advice of 
Pierre Schapira, Robin Hartshorne, Rodney Sharp, 
Manuel Saorin, Suresh Nayak, Liran Shaul, Johan de Jong, Steven Kleiman, 
Louis Rowen, Amnon Neeman, Amiram Braun, Jesse Wolfson, 
Bjarn de Jong, James Zhang, Jun-Ichi Miyachi, Vladimir Hinich, 
Brooke Shipley, Jeremy Rickard, Peter J{\o}rgensen, 
Michael Sharpe (on LaTeX), Roland Berger, Mattia Ornaghi, Georges Maltsiniotis, 
Harry Gindi, Goncalo Tabuada, Sefi Ladkani, Saurabh Singh,
and the anonymous referees for Cambridge University Press. 

Special thanks to Thomas Harris, my editor at Cambridge University Press,  
whose initiative made the publication of this book possible; and to 
Clare Dennison and Libby Haynes at CUP, for their assistance.
